\chapter{The Linux Binary Compatibility Layer}
\label{chap:linux_compat}

Historically, experimental and educational microkernels suffered from an insurmountable paradox: while their internal architectures were elegant, modular, and mathematically verifiable, they were incapable of executing standard desktop software. Authoring an operating system from scratch often meant being confined to a rudimentary text console, running a toy text editor and a handful of basic shell utilities.

Porting complex, multi-million-line applications---such as the **KDE Plasma 6** desktop environment, the **KWin Wayland** compositor, the **Brave/Chromium** web browser, Python 3, and Qt 6---by hand is an engineering impossibility for a small research or academic team.

To bridge this chasm, Hydra introduces a high-performance **Linux Binary Compatibility Subsystem** (\texttt{src/compat/linux/}). Hydra does not run Linux inside a virtual machine or container; instead, it executes unmodified upstream Linux x86\_64 ELF binaries directly on the bare-metal microkernel by providing an in-kernel emulation of the Linux Application Binary Interface (ABI).

In this chapter, we explore how Hydra implements hardware-accelerated system call dispatching, POSIX process and credential emulation, dynamic ELF linking, virtual filesystems (\texttt{/proc}, \texttt{/sys}, \texttt{/dev}), the Wayland graphical stack, and the infrastructure required to host modern web browsers.

% ==============================================================================
\section{The Binary Compatibility Imperative}
\label{sec:compat_imperative}
% ==============================================================================

In the history of operating systems, numerous technically superior systems (such as Plan~9 from Bell Labs, Amoeba, and QNX) struggled to gain widespread commercial adoption because they required applications to be rewritten or recompiled for proprietary APIs. The modern computing ecosystem is dominated by a vast library of open-source software written to POSIX and Linux-specific ABI specifications.

Operating systems architects have historically approached compatibility through three techniques:
\begin{enumerate}
    \item \textbf{Hardware Virtualization (Type-1 or Type-2 Hypervisors):} Running a full guest Linux kernel inside a hardware virtual machine (e.g., WSL~2 on Windows or KVM on Linux). While compatibility is 100\%, this incurs heavy memory overhead (running two kernels simultaneously) and poor graphical composition performance.
    \item \textbf{User-Space Dynamic Translation (e.g., WINE):} Emulating system APIs entirely in user space by intercepting dynamic library symbols. While effective for stable APIs like Win32, it fails when programs make direct inline \texttt{syscall} assembly instructions.
    \item \textbf{Kernel In-Place ABI Emulation (e.g., FreeBSD Linuxulator, WSL~1, and Hydra):} The host kernel natively intercepts the CPU's \texttt{syscall} instruction, evaluates the Linux system call number and arguments, and executes the requested operation directly inside the host kernel with \textbf{zero virtualization overhead}.
\end{enumerate}

Hydra implements this third approach, combining the modularity and safety of a microkernel with the vast software library of Linux.

% ==============================================================================
\section{Hardware System Call Dispatching: \texttt{SYSCALL} and \texttt{SYSRET}}
\label{sec:syscall_dispatch}
% ==============================================================================

On the x86\_64 architecture, user-space applications invoke operating system services using the high-speed \texttt{syscall} instruction rather than legacy software interrupts (such as \texttt{int 0x80}).

During early initialization, Hydra configures the CPU's Model-Specific Registers (MSRs) in \texttt{src/arch/syscall.rs}:
\begin{itemize}
    \item \textbf{\texttt{IA32\_STAR} (MSR \texttt{0xC000\_0081}):} Encodes the 16-bit GDT segment selectors for Ring~0 kernel code/data and Ring~3 user code/data.
    \item \textbf{\texttt{IA32\_LSTAR} (MSR \texttt{0xC000\_0082}):} Holds the 64-bit virtual memory entry point for the microkernel system call handler (\texttt{hydra\_syscall\_dispatcher}).
    \item \textbf{\texttt{IA32\_FMASK} (MSR \texttt{0xC000\_0084}):} Defines the bitmask of flags cleared in \texttt{RFLAGS} upon entering the kernel (e.g., clearing the Interrupt Flag \texttt{IF} to disable interrupts during register preservation).
\end{itemize}

\subsection{The System V AMD64 Calling Convention}
When an application invokes \texttt{syscall}, the processor and kernel adhere to the strict register convention:
\begin{table}[htbp]
\centering
\small
\begin{tabular}{lll}
\hline
\textbf{Register} & \textbf{Role in Syscall Entry} & \textbf{Role in Syscall Return} \\
\hline
\texttt{RAX} & System Call Number (e.g., 0=read, 1=write) & Return value ($\ge 0$ success, $-1$ to $-4095$ errno) \\
\texttt{RDI} & Argument 1 (e.g., file descriptor \texttt{fd}) & Preserved \\
\texttt{RSI} & Argument 2 (e.g., buffer pointer \texttt{*buf}) & Preserved \\
\texttt{RDX} & Argument 3 (e.g., byte count \texttt{count}) & Preserved \\
\texttt{R10} & Argument 4 (e.g., flags or offset) & Preserved \\
\texttt{R8} & Argument 5 & Preserved \\
\texttt{R9} & Argument 6 & Preserved \\
\texttt{RCX} & Overwritten by CPU with Return \texttt{RIP} & Used by \texttt{sysretq} to resume execution \\
\texttt{R11} & Overwritten by CPU with Return \texttt{RFLAGS} & Used by \texttt{sysretq} to restore flags \\
\hline
\end{tabular}
\caption{System V AMD64 System Call Register Convention.}
\label{tab:syscall_regs}
\end{table}

\begin{figure}[htbp]
\centering
\begin{tikzpicture}[scale=0.9, every node/.style={transform shape}]
    % User space
    \node[draw=rustkeyword, fill=white, rounded corners, minimum width=4cm, minimum height=1.2cm] (app) at (0, 3) {
        \begin{tabular}{c}
            \textbf{User Application}\\
            \footnotesize (e.g. \texttt{brave}, \texttt{kwin}, \texttt{python})\\
            \scriptsize Executes \texttt{syscall} (\texttt{RAX=0x01}, \texttt{RDI=1})
        \end{tabular}
    };

    % Hardware transition
    \node[draw=hydragold, fill=yellow!15, line width=1pt, minimum width=3.5cm, minimum height=1.0cm] (hw) at (5.5, 3) {
        \begin{tabular}{c}
            \textbf{x86\_64 CPU Hardware}\\
            \scriptsize MSR \texttt{LSTAR} Transition\\
            \scriptsize Ring 3 $\rightarrow$ Ring 0
        \end{tabular}
    };

    % Kernel dispatcher
    \node[draw=hydrablue, fill=hydralight, rounded corners, minimum width=4cm, minimum height=1.4cm] (dispatch) at (11, 3) {
        \begin{tabular}{c}
            \textbf{Hydra Syscall Dispatcher}\\
            \footnotesize \texttt{src/compat/linux/mod.rs}\\
            \scriptsize Demux \texttt{RAX}: \texttt{SYS\_write(1, buf, len)}
        \end{tabular}
    };

    \draw[->, >=Stealth, line width=1.5pt, rustkeyword] (app) -- (hw);
    \draw[->, >=Stealth, line width=1.5pt, hydragold] (hw) -- (dispatch);
    \draw[->, >=Stealth, dashed, line width=1.2pt, rustred] (dispatch) to[bend left=30] node[below] {\footnotesize \texttt{sysretq} (Result in \texttt{RAX})} (app);
\end{tikzpicture}
\caption{The x86\_64 \texttt{syscall} / \texttt{sysret} Dispatch Pipeline in Hydra.}
\label{fig:syscall_pipeline}
\end{figure}

The assembly dispatcher saves the caller's registers onto the thread's kernel stack, calls into the Rust dispatcher \texttt{dispatch\_linux\_syscall()} in \texttt{src/compat/linux/mod.rs}, and deposits the result into \texttt{RAX}. Returning to user space requires a single execution of \texttt{sysretq}, re-entering Ring~3 with near zero cycle overhead.

% ==============================================================================
\section{POSIX Process Model and Thread-Local Storage (TLS)}
\label{sec:posix_model}
% ==============================================================================

Modern software frameworks (glibc, musl, Qt 6, LLVM) rely on sophisticated POSIX process semantics:

\subsection{Process Creation via \texttt{clone}}
In Linux, \texttt{fork()}, \texttt{vfork()}, and \texttt{pthread\_create()} are all thin wrappers around the versatile \texttt{clone(2)} system call. Hydra provides complete support for \texttt{clone}:
\begin{itemize}
    \item \texttt{CLONE\_VM}: The child shares the parent's virtual address space (creating a thread).
    \item \texttt{CLONE\_FS}: The child shares file system information (working directory, root directory, umask).
    \item \texttt{CLONE\_FILES}: The child shares the file descriptor table.
    \item \texttt{CLONE\_SIGHAND}: The child shares signal handlers.
    \item \texttt{CLONE\_THREAD}: The child is placed in the same thread group as the parent.
    \item \texttt{CLONE\_SETTLS}: Configures Thread-Local Storage for the new thread.
    \item \texttt{CLONE\_CHILD\_CLEARTID}: Automatically clears a futex word at the child's TID address and wakes waiting parent threads upon termination.
\end{itemize}

\subsection{Thread-Local Storage and \texttt{arch\_prctl}}
In multithreaded programs, variables declared with \texttt{thread\_local} or \texttt{\_\_thread} must have a unique instance per thread. On x86\_64, compilers address thread-local variables using offsets from the \texttt{FS} segment register (e.g., \texttt{mov \%fs:0x10, \%rax}).

Because segment registers cannot hold 64-bit base addresses directly, the processor provides the **\texttt{IA32\_FS\_BASE} MSR** (\texttt{0xC000\_0100}). When an application calls \texttt{arch\_prctl(ARCH\_SET\_FS, addr)}, Hydra writes the 64-bit base address directly into \texttt{IA32\_FS\_BASE} and saves it in the thread's \texttt{TaskControlBlock}. During every context switch, the scheduler automatically restores the incoming thread's \texttt{fs\_base}.

\subsection{Signal Handling and User Stack Trampolines}
Signals provide asynchronous event notification in POSIX systems (e.g., \texttt{SIGINT}, \texttt{SIGSEGV}, \texttt{SIGCHLD}). Delivering a signal to a user thread requires careful stack manipulation:
\begin{enumerate}
    \item The kernel builds a **Signal Frame** directly on the user thread's stack, pushing the saved register context (\texttt{ucontext\_t}) and signal information (\texttt{siginfo\_t}).
    \item The kernel rewrites the thread's \texttt{RIP} to point to the user-registered signal handler function.
    \item The kernel sets the thread's return address to point to a **Signal Trampoline** containing the assembly instruction \texttt{syscall} with \texttt{RAX = SYS\_rt\_sigreturn}.
    \item When the handler finishes, it returns to the trampoline, which invokes \texttt{rt\_sigreturn}. The kernel restores the original register context from the stack frame and resumes ordinary execution without glitching.
\end{enumerate}

% ==============================================================================
\section{Virtual Filesystems: \texttt{/proc}, \texttt{/sys}, and \texttt{/dev}}
\label{sec:virtual_filesystems}
% ==============================================================================

Modern desktop applications, dynamic linkers, and system libraries do not interact exclusively with regular files stored on physical media; they inspect synthetic virtual filesystems generated dynamically in memory by the kernel:

\begin{table}[htbp]
\centering
\small
\begin{tabular}{lll}
\hline
\textbf{Virtual Path} & \textbf{Type} & \textbf{Synthesized Contents and Purpose} \\
\hline
\texttt{/proc/self/cmdline} & Text & Null-delimited command-line arguments of running process \\
\texttt{/proc/self/exe} & Symlink & Canonical path to the running ELF binary on the ZFS root \\
\texttt{/proc/self/maps} & Text & Virtual memory regions, permissions (rwxp), and mapped files \\
\texttt{/proc/cpuinfo} & Text & Processor vendor, model, clock speed, and feature flags \\
\texttt{/proc/meminfo} & Text & Total RAM, free frames, buffer cache, and ZFS pool stats \\
\texttt{/sys/class/drm/} & Directory & DRM display cards, connectors, and framebuffers \\
\texttt{/dev/null} & Char Dev & Discards all writes, returns EOF on read (Major 1, Minor 3) \\
\texttt{/dev/zero} & Char Dev & Streams infinite null bytes (\texttt{0x00}) for memory clearing \\
\texttt{/dev/urandom} & Char Dev & Cryptographically secure pseudo-random bytes from hardware \\
\texttt{/dev/dri/card0} & Char Dev & Direct Rendering Manager node for display compositing \\
\texttt{/dev/shm} & Tmpfs & POSIX Shared Memory with sticky-bit permissions (\texttt{01777}) \\
\hline
\end{tabular}
\caption{Key Virtual Filesystem Nodes Implemented in Hydra.}
\label{tab:vfs_nodes}
\end{table}

In Hydra, these files are synthesized in memory on demand through \texttt{src/compat/linux/proc.rs} and \texttt{src/compat/linux/vfs.rs}. When an application reads \texttt{/proc/self/maps}, the kernel traverses the task's active memory mappings and formats human-readable text into the caller's read buffer.

% ==============================================================================
\section{Hosting Modern Desktop Environments (SDDM, KWin, KDE Plasma 6)}
\label{sec:hosting_plasma}
% ==============================================================================

The ultimate verification of an operating system's architectural completeness is its ability to host complex, graphical user applications. Hydra boots directly into a complete **KDE Plasma 6** desktop:

\begin{figure}[htbp]
\centering
\begin{tikzpicture}[scale=0.9, every node/.style={transform shape}]
    % Layer 4: Desktop apps
    \node[draw=rustkeyword, fill=white, rounded corners, minimum width=11cm, minimum height=0.8cm] (apps) at (5.5, 4.2) {
        \textbf{Applications:} Brave Web Browser, Konsole Terminal, Dolphin File Manager, Python
    };

    % Layer 3: Compositor and Display Manager
    \node[draw=hydrablue, fill=hydralight, rounded corners, minimum width=11cm, minimum height=0.8cm] (gui) at (5.5, 2.8) {
        \textbf{Wayland Graphics Stack:} SDDM Greeter, KWin Wayland Compositor, PipeWire
    };

    % Layer 2: System Daemons
    \node[draw=hydragold, fill=white, rounded corners, minimum width=11cm, minimum height=0.8cm] (daemons) at (5.5, 1.4) {
        \textbf{System Services:} D-Bus System Bus, \texttt{hydra-logind} (Seat/Session Manager), Polkit
    };

    % Layer 1: Hydra Microkernel + Linux ABI
    \node[draw=rustred, fill=rustred!10, line width=1.2pt, rounded corners, minimum width=11cm, minimum height=0.8cm] (core) at (5.5, 0.0) {
        \textbf{Hydra Microkernel Core:} Linux ABI Dispatcher, OpenZFS Engine, Isolated Ring 3 Drivers
    };

    \draw[->, >=Stealth, line width=1.2pt] (apps) -- (gui);
    \draw[->, >=Stealth, line width=1.2pt] (gui) -- (daemons);
    \draw[->, >=Stealth, line width=1.2pt] (daemons) -- (core);
\end{tikzpicture}
\caption{The Full Application Stack Running on Top of Hydra.}
\label{fig:plasma_stack}
\end{figure}

The graphical bootstrap pipeline proceeds through five orchestrated phases:
\begin{enumerate}
    \item \textbf{System Message Bus:} The boot supervisor launches \texttt{dbus-daemon --system}, providing interprocess communication for desktop components over UNIX domain sockets (\texttt{/var/run/dbus/system\_bus\_socket}).
    \item \textbf{Session and Seat Management (\texttt{hydra-logind}):} Modern desktop environments require an active seat manager. Hydra launches \texttt{hydra-logind} (\texttt{userland/hydra-logind/}), which connects to D-Bus, implements the \texttt{org.freedesktop.login1} interface, registers \texttt{seat0}, and activates the graphical virtual terminal.
    \item \textbf{Display Manager (SDDM):} The Simple Desktop Display Manager starts, loads its configuration from \texttt{/usr/share/sddm/themes/breeze/theme.conf}, initializes PAM authentication, and displays the graphical login greeter.
    \item \textbf{The KWin Wayland Compositor:} Upon user authentication, SDDM spawns \texttt{kwin\_wayland}. KWin opens \texttt{/dev/dri/card0}, initializes DRM dumb framebuffers, creates the Wayland display socket (\texttt{/run/user/1000/wayland-0}), and launches \texttt{plasmashell}.
    \item \textbf{Full Application Execution:} Applications (such as Konsole, Dolphin, and Brave) connect to the Wayland socket, submit drawing surfaces via shared memory (\texttt{/dev/shm}), and receive input events from \texttt{/dev/input/event0} and \texttt{event1}.
\end{enumerate}

% ==============================================================================
\section{Hosting Modern Web Browsers: Brave and Chromium}
\label{sec:hosting_brave}
% ==============================================================================

A modern web browser is effectively an operating system in its own right, incorporating JavaScript JIT compilers, WebAssembly runtimes, multithreaded rendering engines, hardware video decoders, and sandboxed security containers.

To host the **Brave Browser** successfully on Hydra, three critical kernel facilities are provided:
\begin{enumerate}
    \item \textbf{POSIX Shared Memory (\texttt{/dev/shm}):} Chromium architectures partition tasks across independent processes: the main Browser process, sandboxed Renderer processes, and the GPU process. High-speed rendering occurs by rendering pixel buffers into shared memory frames created via \texttt{memfd\_create()} or \texttt{shm\_open()} on \texttt{/dev/shm}.
    \item \textbf{Cryptographic NSS Databases:} Browsers maintain cryptographic certificates and security keys in an on-disk database (\texttt{sql:\textasciitilde/.pki/nssdb}). Hydra provides complete POSIX file locking (\texttt{fcntl} record locking via \texttt{F\_SETLK} and \texttt{F\_GETLK}) to ensure atomic SQLite3 database transactions without corruption.
    \item \textbf{Asynchronous Epoll Networking:} The browser performs concurrent DNS lookups and TLS handshakes over non-blocking sockets, multiplexed through Hydra's $O(1)$ \texttt{epoll} implementation.
\end{enumerate}

% ==============================================================================
\section{Summary}
\label{sec:summary_chap7}
% ==============================================================================

In this chapter, we explored how Hydra resolves the historical dilemma of microkernel software isolation:
\begin{itemize}
    \item Direct bare-metal ABI emulation avoids the severe performance penalties of virtual machines while providing 100\% binary compatibility.
    \item Fast \texttt{syscall}/\texttt{sysret} MSR configuration delivers microsecond system call dispatching.
    \item Comprehensive POSIX emulation supports multithreaded process trees (\texttt{clone}), Thread-Local Storage (\texttt{FS\_BASE}), and signal delivery frames.
    \item Dynamic virtual filesystems (\texttt{/proc}, \texttt{/sys}, \texttt{/dev}, \texttt{/dev/shm}) provide standard Linux kernel introspection and memory sharing.
    \item Hydra successfully runs complete, unmodified upstream software stacks, from D-Bus and SDDM to KWin Wayland, KDE Plasma 6, and modern Chromium-based web browsers.
\end{itemize}

% ==============================================================================
\section{Exercises and Review Problems}
\label{sec:exercises_chap7}
% ==============================================================================

\begin{enumerate}
    \item \textbf{ABI Compatibility vs. API Compatibility:} What is the difference between Application Binary Interface (ABI) compatibility and Application Programming Interface (API) compatibility? Why is ABI compatibility necessary to execute pre-compiled Linux binaries?
    \item \textbf{Syscall Register Conventions:} List the x86\_64 CPU registers used to pass arguments in the Linux system call convention. How does the kernel signal an error return value to the caller?
    \item \textbf{Fast Syscall Mechanics:} Explain the operation of the \texttt{IA32\_LSTAR} and \texttt{IA32\_STAR} MSRs. Why is the \texttt{syscall} instruction faster than the traditional \texttt{int 0x80} software interrupt?
    \item \textbf{Thread-Local Storage:} What is the purpose of the \texttt{IA32\_FS\_BASE} register? How does Hydra update this register when switching between threads?
    \item \textbf{Signal Trampolines:} Why does the kernel build a stack frame on the user thread's stack when delivering a signal? What is the function of the \texttt{rt\_sigreturn} system call?
    \item \textbf{Virtual Filesystem Introspection:} When an application opens and reads \texttt{/proc/self/maps}, where does the data originate? Why does this file consume zero bytes on physical storage media?
    \item \textbf{Shared Memory Permissions:} Why must the \texttt{/dev/shm} directory be configured with permissions \texttt{01777} (world-writable with the sticky bit set)? What security vulnerability arises if the sticky bit is omitted?
    \item \textbf{Desktop Session Architecture:} Explain the role of D-Bus and the seat manager (\texttt{hydra-logind}) in starting a graphical desktop environment. Why does KWin Wayland require an active logind session before opening \texttt{/dev/dri/card0}?
    \item \textbf{Browser Multi-Process Sandboxing:} Describe how modern browsers use separate processes for rendering web pages. How does Hydra's \texttt{/dev/shm} enable zero-copy graphical composition between the renderer and the browser UI?
    \item \textbf{Design Problem:} Design an emulation handler for the Linux \texttt{futex(2)} system call with operations \texttt{FUTEX\_WAIT} and \texttt{FUTEX\_WAKE}. How does your design ensure atomicity between inspecting the user memory value and enqueuing the calling thread on the wait queue?
\end{enumerate}
